\section*{Разбор задачи из тестового варианта №1}
\addcontentsline{toc}{section}{Разбор задачи из тестового варианта №1}
\begin{Task}{Асимптотическое поведение экстремумов}{task1}
Случайные величины $\{\xi_n, n \in \mathbb{N}\}$ независимы и имеют распределение $\Gamma(2, 1)$ с плотностью
\[
p(x) = 4xe^{-2x}\I\{x > 0\}.
\]
Найдите с вероятностью 1 значение
\[
\varlimsup_{n \to \infty} \frac{\xi_n - \frac{1}{2}\ln n}{\ln \ln n}.
\]
\end{Task}

\begin{Solution}
Требуется найти константу $C$, для которой $\P\left(\varlimsup_{n\to\infty} \eta_n = C\right) = 1$, где 
\[
\eta_n = \frac{\xi_n - \frac{1}{2}\ln n}{\ln \ln n}.
\]
Математическим аппаратом для поиска почти верных пределов последовательностей независимых случайных величин служат леммы Бореля\,---\,Кантелли (ЛБК). 

\begin{MethodBox}[Анализ верхнего предела последовательности]
Для доказательства того, что $\varlimsup_{n\to\infty} \eta_n = C$ почти наверное, необходимо для произвольного $\varepsilon > 0$ исследовать сходимость рядов вероятностей:
\begin{enumerate}
    \item \textbf{Оценка сверху:} Если ряд $\sum_{n} \P(\eta_n > C + \varepsilon)$ сходится, то по первой лемме Бореля\,---\,Кантелли событие происходит конечное число раз, следовательно, $\varlimsup \eta_n \leqslant C$ п.н.
    \item \textbf{Оценка снизу:} Если ряд $\sum_{n} \P(\eta_n > C - \varepsilon)$ расходится, то (при условии независимости событий) по второй лемме Бореля\,---\,Кантелли событие происходит бесконечно часто, следовательно, $\varlimsup \eta_n \geqslant C$ п.н.
\end{enumerate}
\end{MethodBox}

\textbf{Шаг 1. Асимптотика хвоста распределения.} \\
Вычислим вероятность того, что случайная величина $\xi_n$ превышает некоторый порог $x > 0$. Интегрируя по частям, получаем:
\begin{align*}
\P(\xi_n > x) &= \int_{x}^{+\infty} 4t e^{-2t} \, dt = \left[ \begin{array}{ll} 
u = 2t & du = 2\,dt \\ 
dv = 2e^{-2t}\,dt & v = -e^{-2t} \end{array} \right] = \\
&= \left. -2t e^{-2t} \right|_{x}^{+\infty} + \int_{x}^{+\infty} 2e^{-2t} \, dt = 2xe^{-2x} + e^{-2x} = (2x + 1)e^{-2x}.
\end{align*}
При $x \to +\infty$ старшим членом является $2x$, поэтому справедлива асимптотическая эквивалентность:
\[
\P(\xi_n > x) \sim 2x e^{-2x}.
\]

\textbf{Шаг 2. Оценка вероятностей отклонений.} \\
Рассмотрим событие $A_n(C) = \{\eta_n > C\}$. Выразим $\xi_n$ из неравенства:
\[
\P(\eta_n > C) = \P\left(\xi_n > \frac{1}{2}\ln n + C \ln \ln n\right).
\]
Обозначим критическую границу как $x_n = \frac{1}{2}\ln n + C \ln \ln n$. Заметим, что при $n \to \infty$ граница $x_n \to +\infty$, что позволяет использовать найденную асимптотику хвоста:
\[
\P(A_n(C)) \sim 2x_n e^{-2x_n}.
\]
Вычислим асимптотическое поведение каждого множителя по отдельности:
\begin{align*}
2x_n &= \ln n + 2C \ln \ln n \sim \ln n, \\
e^{-2x_n} &= \exp\left(-\ln n - 2C \ln \ln n\right) = \frac{1}{n (\ln n)^{2C}}.
\end{align*}
Перемножая полученные выражения, находим асимптотику вероятности:
\[
\P(A_n(C)) \sim (\ln n) \cdot \frac{1}{n (\ln n)^{2C}} = \frac{1}{n (\ln n)^{2C - 1}}.
\]

\textbf{Шаг 3. Применение лемм Бореля\,---\,Кантелли.} \\
Полученная вероятность представляет собой общий член обобщенного гармонического ряда (ряда Бертрана) $\sum \frac{1}{n (\ln n)^p}$, который сходится при $p > 1$ и расходится при $p \leqslant 1$. В нашем случае параметр степени $p = 2C - 1$.

Рассмотрим два случая для произвольного $\varepsilon > 0$:
\begin{enumerate}
    \item Пусть $C = 1 + \varepsilon$. Тогда степень логарифма равна $2(1 + \varepsilon) - 1 = 1 + 2\varepsilon > 1$.
    Ряд $\sum \P(\eta_n > 1 + \varepsilon)$ сходится. По I~лемме Бореля\,---\,Кантелли, с вероятностью 1 реализуется лишь конечное число событий $A_n(1+\varepsilon)$. Отсюда:
    \[
    \varlimsup_{n \to \infty} \eta_n \leqslant 1 + \varepsilon \quad \text{п.н.}
    \]
    В силу произвольности $\varepsilon$, получаем $\varlimsup_{n \to \infty} \eta_n \leqslant 1$.

    \item Пусть $C = 1 - \varepsilon$. Тогда степень логарифма равна $2(1 - \varepsilon) - 1 = 1 - 2\varepsilon < 1$.
    Ряд $\sum \P(\eta_n > 1 - \varepsilon)$ расходится. Поскольку исходные случайные величины $\xi_n$ независимы, события $A_n(1-\varepsilon)$ также независимы. По II~лемме Бореля\,---\,Кантелли, с вероятностью 1 реализуется бесконечно много таких событий. Отсюда:
    \[
    \varlimsup_{n \to \infty} \eta_n \geqslant 1 - \varepsilon \quad \text{п.н.}
    \]
    В силу произвольности $\varepsilon$, получаем $\varlimsup_{n \to \infty} \eta_n \geqslant 1$.
\end{enumerate}
Объединяя оба результата, делаем вывод, что верхний предел в точности равен 1 почти наверное.
\end{Solution}

\begin{Answer}
1
\end{Answer}
